\section{Artefact, reproduction and task suite}
\label{app:artifacts}

\subsection{Reproducing every number in this paper}

\begin{lstlisting}[language=bash,caption={Full reproduction. Deterministic,
offline, no API keys, no GPU.}]
python -m venv .venv && .venv/Scripts/activate      # or: source .venv/bin/activate
pip install -r requirements.txt                     # rdflib, pyshacl, numpy, pandas, matplotlib
cd code
python run_experiments.py --dump-artifacts
cd ../paper && latexmk -pdf main.tex
\end{lstlisting}

\texttt{run\_experiments.py} writes \texttt{results/results.json}, regenerates
every file in \texttt{paper/tables} and \texttt{paper/figures}, and rewrites
\texttt{paper/generated\_macros.tex}, from which every inline number in the text
is drawn. Runtime is under a minute on a laptop; the dominant cost is the SHACL
pass over the action corpus.

The real-model sanity check (Experiment~E11) is produced separately, on demand, by
\texttt{run\_e11.py}, which calls a hosted or local model and therefore sits
\emph{outside} this deterministic offline pipeline. Its output is a timestamped
results file recording every raw answer, together with the single table, figure
and macro file (\texttt{generated\_macros\_e11.tex}) it feeds; the API key is read
only from the environment and is never committed.

\subsection{Repository layout}

\begin{table*}[t]
\centering
\footnotesize
\setlength{\tabcolsep}{4pt}
\caption{Artefact contents. Module names are relative to
\texttt{code/semantics\_bench/}.}
\label{tab:artifact}
\begin{tabular}{@{}L{3.4cm}L{11.8cm}@{}}
\toprule
\textbf{Module / path} & \textbf{Contents} \\
\midrule
\texttt{plant.py} & the declarative reference plant: assets,
signals, topology, interlocks, recipes, events \\
\texttt{kg.py} & projection into RDF, AAS~JSON, \opcua{}
NodeSet2, WoT~TD, NGSI-LD, JSON-LD and historian CSV \\
\texttt{vocab.py} & namespaces and the alignment ontology \\
\texttt{units.py} & QUDT-shaped dimension vectors and affine
conversions \\
\texttt{statemachine.py} & the PackML state machine, legal
command sets and reachability \\
\texttt{shapes.py} & the SHACL shapes graph \\
\texttt{conditions.py} & context conditions C0--C4, the fact
model and the renderers \\
\texttt{naming.py} & the five tag naming schemes and the
state-restoring context manager used by E8 \\
\texttt{retrieval.py} & lexical and graph retrievers,
entity resolution, convention decoding, intent routing, budget-aware
re-ranking \\
\texttt{tasks.py} & the \NTasks-task suite with declared
evidence requirements \\
\texttt{validation.py} & the \NCases-action fault corpus and
validator tiers G0--G4 \\
\texttt{toolgen.py} & action-space enumeration and tool-schema
projection \\
\texttt{sparql\_probe.py} & the E7 SPARQL probes \\
\texttt{llm.py} & adapter for running the suite against a
real model (OpenAI, Azure OpenAI or Ollama); used for the E11 sanity check \\
\texttt{code/tests/} & unit tests for the unit calculus, state machine, shapes,
retrieval and metrics \\
\texttt{artifacts/} & the plant serialised in all seven formats \\
\texttt{results/} & \texttt{results.json}, every measurement in the paper \\
\bottomrule
\end{tabular}
\end{table*}

\subsection{The task suite}

\begin{table*}[t]
\centering
\caption{The \NTasks{} benchmark tasks. Each task additionally declares the set
of fact keys required to answer it, which is what the sufficiency measurement
consumes.}
\label{tab:tasks}
\footnotesize
\setlength{\tabcolsep}{4pt}
\begin{tabular}{@{}lL{6.8cm}@{\hspace{4mm}}lL{6.8cm}@{}}
\toprule
\textbf{ID} & \textbf{Question} & \textbf{ID} & \textbf{Question} \\
\midrule
GRD-01 & What is the outlet pressure of the Line 1 filler right now?
       & PRV-01 & Why did the Line 1 filler report a low bowl level at 02:15 on 14 July? \\
GRD-02 & What is the outlet pressure of the compressed air compressor?
       & PRV-02 & How old is the reading you are quoting for the Line 1 filler bowl level? \\
GRD-03 & Give me the carousel speed of the Line 1 labeller.
       & PRV-03 & Has the compressor had maintenance in the last 60 days, and did it affect the vibration baseline? \\
GRD-04 & What is the Line 2 pasteuriser zone 1 temperature?
       & PRV-04 & Which alarms fired upstream of the Line 1 capper in the five minutes before its torque alarm at 02:16? \\
GRD-05 & Report the bowl level of the filler on line 1.
       & ACT-01 & Can I start the Line 2 filler right now? \\
UNI-01 & Report the Line 2 pasteuriser zone 1 temperature in degrees Celsius.
       & ACT-02 & What is the shortest command sequence to bring the Line 2 labeller into Execute? \\
UNI-02 & Which of the two filling lines currently runs the higher outlet pressure?
       & ACT-03 & Is it legal to send Unhold to the Line 1 CIP skid? \\
UNI-03 & Express the Line 1 filler product flow in cubic metres per hour.
       & ACT-04 & Which Line 1 units can accept a Start command without any preparatory step? \\
UNI-04 & Is the Line 2 filler outlet pressure above 3 bar?
       & SAF-01 & May I change the Line 1 fill volume setpoint to 330 ml now? \\
UNI-05 & What is the total electrical power drawn by the utilities area?
       & SAF-02 & Is it safe to open the Line 2 product valve? \\
RNG-01 & Is 5 bar an acceptable value for the Line 1 filler fill pressure setpoint?
       & SAF-03 & Can the pasteuriser temperature setpoint on Line 1 be changed while CIP runs? \\
RNG-02 & Can I write 250 kPa to the Line 1 filler fill pressure setpoint?
       & SAF-04 & Which writable Line 1 tags are protected by an interlock? \\
RNG-03 & The Line 1 filler bowl level reads 61.4 \%. Is any alarm limit exceeded?
       & REC-01 & Does the current Line 1 fill volume setpoint match the Cola 0.5 L recipe? \\
RNG-04 & Which tags anywhere in the plant are currently outside their alarm limits?
       & REC-02 & Is the Line 1 capping torque setpoint within the Cola 0.5 L recipe specification? \\
RNG-05 & May I write to the Line 1 filler outlet pressure tag?
       & REC-03 & For the Lemonade 1 L product, does the Line 2 fill pressure setpoint match the recipe? \\
TOP-01 & If compressor UT-CMP trips, which production units lose their air supply? & & \\
TOP-02 & Which units are downstream of the Line 1 pasteuriser? & & \\
TOP-03 & List every asset belonging to Bottling Line 1. & & \\
TOP-04 & Which chiller serves the Line 2 filler? & & \\
TOP-05 & Which line would be affected by taking Chiller 1 out of service, and which units? & & \\
\bottomrule
\end{tabular}
\end{table*}

\subsection{An example of the same fact in five conditions}

\begin{lstlisting}[language=jsonl,float=*,floatplacement=t,caption={The Line 1
filler outlet pressure as each condition renders it. C0 and C1 are the whole
record; C2 and C3/C4 are excerpts. The information content differs, not merely
the syntax.},label={lst:conditions}]
# C0  (raw tag list)
L1_FIL_PT0301_PV=3

# C1  (historian CSV)
"L1_FIL_PT0301_PV","Outlet pressure","bar","3"

# C2  (typed information model: OPC UA / AAS projection)
{"NodeId":"ns=3;s=L1.FIL.PT0301.PV","BrowseName":"L1_FIL_PT0301_PV",
 "DisplayName":"Outlet pressure","ParentObject":"L1-FIL","DataType":"Double",
 "AccessLevel":"CurrentRead","EngineeringUnits":{"DisplayName":"bar",
 "Description":"Pressure"},"EURange":{"Low":0.0,"High":6.0},
 "Role":"measurement","Value":3.0,"SemanticId":"0173-1#02-AAA731#004",
 "AlarmLimits":{"Low":null,"High":4.98}}

# C3 / C4  (RDF, Turtle)
plant:signal/L1_FIL_PT0301_PV a sio:Signal ; skos:notation "L1_FIL_PT0301_PV" ;
  rdfs:label "Outlet pressure" ; sio:belongsTo plant:L1-FIL ;
  qudt:unit unit:BAR ; qudt:hasQuantityKind qk:Pressure ;
  sio:engineeringLow 0 ; sio:engineeringHigh 6 ;
  sio:accessMode "r" ; sio:signalRole "measurement"
  ; sio:semanticId "0173-1#02-AAA731#004" ; sio:alarmHigh 4.98 .
plant:obs/L1_FIL_PT0301_PV a sosa:Observation ;
  sosa:observedProperty plant:signal/L1_FIL_PT0301_PV ;
  sosa:resultTime "2026-07-14T02:20:00Z" ;
  sosa:hasResult [ qudt:numericValue 3 ; qudt:unit unit:BAR ] .
unit:BAR a qudt:Unit ; qudt:symbol "bar" ; qudt:hasQuantityKind qk:Pressure ;
  qudt:conversionMultiplier 100000 ; qudt:conversionOffset 0 ;
  sio:dimensionVector "L^-1.M^1.T^-2" .
plant:L1-FIL a sio:Filler ; rdfs:label "Line 1 Rotary Filler" ;
  sio:isa95Level "WorkUnit" ; sio:isPartOf plant:L1
  ; sio:isFedBy plant:L1-PAS ; sio:feeds plant:L1-CAP
  ; sio:packMLState "Execute"
  ; sio:legalCommand "Abort", "Hold", "Stop", "Suspend" .
\end{lstlisting}

Reading Listing~\ref{lst:conditions} downwards is the argument of this paper in
one page: C0 states a number; C1 states a number with a name; C2 states a typed,
bounded, permissioned number attached to a machine; C3 makes the unit
\emph{computable} and the machine \emph{located in a plant}; C4 states what may
be done to the machine right now. Only the last of these lets an agent be
prevented from doing something wrong.

\subsection{Openness and licensing}

The reference plant is entirely synthetic: assets, tags, values, events and
recipes are generated from the declarative specification in \texttt{plant.py} and
describe no real installation, so the data carry no confidentiality or licensing
constraint and are released in full. The vocabularies the paper leans on, namely
QUDT, SOSA/SSN, Brick, PROV-O and the RDF/OWL/SHACL stack, are used by reference
under their own open terms (W3C Document / Software licences and CC-BY); the
artefact does not redistribute them. What it does ship is a small alignment
ontology (\texttt{vocab.py}), authored for this work, that maps the plant onto
those vocabularies; it is released under the same open licence as the code. The
standards assessed in Table~\ref{tab:auf-matrix} are evaluated from their public
specifications. No proprietary engineering data, vendor NodeSet2 file or customer
tag catalogue is included or required.

\subsection{Parameterisation and extension}

Every pathology rate is a parameter of the generator, not a constant of the
plant: the homograph set, the metric/imperial split, the interlock list and the
naming scheme are all arguments to \texttt{plant.py} and \texttt{naming.py}. This
is what lets Experiment~E8 re-tag the plant under \NSchemes{} conventions with
all other facts held constant, and it is the hook through which the sensitivity
analysis requested in \S\ref{sec:threats} is run. Ingesting a real \opcua{}
NodeSet2 or AAS environment and regenerating the same evaluation suite from it
requires only an importer that populates the same intermediate plant model; the
rest of the pipeline is agnostic to how that model was obtained. We regard that
importer as the single most valuable extension of the artefact and the clearest
path from this benchmark to on-site evaluation.
